\documentclass[12pt]{article}
\usepackage[margin=1in]{geometry}
\usepackage[utf8]{inputenc}
\usepackage{latexsym, amsfonts, enumitem, graphicx, environ,enumitem, fancyhdr, color, amssymb, amsthm, amsmath, mathtools, tikz, nicematrix}
\usetikzlibrary{patterns}
\usetikzlibrary{decorations.pathreplacing}
\pagestyle{fancy}
\setlength{\headheight}{65pt}
\newenvironment{problem}[2][Problem]{\begin{trivlist}
    \item[\hskip \labelsep {\bfseries #1}\hskip \labelsep {\bfseries #2.}]}{\end{trivlist}}
\DeclareMathOperator{\Ima}{Im}

\usepackage{hyperref}
\usepackage[capitalize]{cleveref}

\newtheorem{thm}{Theorem}[section]
\newtheorem{lem}[thm]{Lemma}
\newtheorem{cl}[thm]{Claim}
\newtheorem{prop}[thm]{Proposition}
\newtheorem{cor}[thm]{Corollary}
\newtheorem{conj}[thm]{Conjecture}
\newtheorem{obstruction}[thm]{Obstruction}

\usepackage{mdframed}

\theoremstyle{definition}
\newtheorem{definition}[thm]{Definition}
\newtheorem{remark}[thm]{Remark}
\newtheorem{question}{Question}
\newtheorem{example}[thm]{Example}
\newtheorem{exercise}[thm]{Exercise}
\newtheorem{properties}[thm]{Properties}

\usepackage{tcolorbox}
\tcbuselibrary{theorems,skins,breakable}

\tcolorboxenvironment{example}{
enhanced jigsaw,colframe=cyan,interior hidden,
breakable,%before skip=10pt,after skip=10pt
}

\tcolorboxenvironment{thm}{
enhanced jigsaw,colframe=red,interior hidden,
breakable,%before skip=10pt,after skip=10pt
}

\tcolorboxenvironment{prop}{
enhanced jigsaw,colframe=red,interior hidden,
breakable,%before skip=10pt,after skip=10pt
}

\tcolorboxenvironment{properties}{
enhanced jigsaw,colframe=red,interior hidden,
breakable,%before skip=10pt,after skip=10pt
}

\tcolorboxenvironment{lem}{
enhanced jigsaw,colframe=red,interior hidden,
breakable,%before skip=10pt,after skip=10pt
}
\tcolorboxenvironment{cor}{
enhanced jigsaw,colframe=red,interior hidden,
breakable,%before skip=10pt,after skip=10pt
}
\tcolorboxenvironment{obstruction}{
enhanced jigsaw,colframe=red,interior hidden,
breakable,%before skip=10pt,after skip=10pt
}

\tcolorboxenvironment{proof}{
enhanced jigsaw, frame hidden,%interior hidden,
breakable,%before skip=10pt,after skip=10pt
}

\lhead{Avi Chetlin \\ Assignment 02 \\ 09 September 2026}
\rhead{Dr. David Singer \\ MATH 465 Differential Geometry \\ Fall 2026}

\begin{document}
Unless explicitly stated, \(\alpha: I \rightarrow \mathbb{R}^3\) is a curve parametrized by arc length, \( s \), with curvature \( k(s) \neq 0 \), for all \( s \in I \).

\begin{problem}{1.5.1}
  Given the parametrized curve (helix)
  \begin{equation}
  \label{eq:1}
  \alpha(s) = \left( a \cos \frac{s}{c}, a \sin \frac{s}{c}, b \frac{s}{c} \right), \quad s \in \mathbb{R},
  \end{equation}
  where \( c^2 = a^2 + b^2 \),

  \begin{enumerate}[label=(\alph*)]
    \item Show that the parameter \( s \) is the arc length.
          \begin{proof}
            To prove this, we show that for arbitrary \( s \in I \), the arc length satisfies
            \begin{equation}
            \label{eq:3}
            \int\limits_0^s \left| \alpha^{\prime}(t) \right| \,dt = s.
            \end{equation}
            First of all,
            \begin{equation}
            \label{eq:4}
            \alpha^{\prime}(s) = \left( -\frac{a}{c}\sin \frac{s}{c}, \frac{a}{c}\cos \frac{s}{c}, \frac{b}{c} \right),
            \end{equation}
            so
            \begin{equation}
            \label{eq:5}
            \left| \alpha^{\prime} \right|^2 = \left( \frac{a}{c} \right)^2 + \left( \frac{b}{c} \right)^2 = \frac{a^2 + b^2}{c^2} = 1.
            \end{equation}
            Thus, our arc length integral becomes
            \begin{equation}
            \label{eq:6}
            \int\limits_0^s dt \equiv s.
            \end{equation}
          \end{proof}
    \item Determine the curvature and the torsion of \(\alpha\).
          \begin{proof}

          \end{proof}
    \item Determine the osculating plane of \(\alpha\).
    \item Show that the lines containing \( n(s) \) and passing through \(\alpha(s)\) meet the \( z \)-axis under a constant angle equal to \( \frac{\pi}{2} \).
    \item Show that the tangent lines to \(\alpha\) make a constant angle with the \( z \)-axis.
  \end{enumerate}
\end{problem}

\begin{problem}{1.5.14}
  Let \(\alpha:(a,b) \rightarrow \mathbb{R}^2\) be a regular parametrized plane curve.
  Assume that there exists \( t_0 \), \( a < t_0 < b \), such that the distance \( \left| \alpha(t) \right| \) from the origin to the trace of \(\alpha\) will be maximum at \( t_{0} \).
  Prove that the curvature \( k \) of \(\alpha\) at \( t_0 \) satisfies \( \left| k(t_0) \right| \geq 1/\left| \alpha(t_0) \right|\).
\end{problem}

\begin{problem}{1.5.17}
  In general, a curve \(\alpha\) is called a \textit{helix} if the tangent lines of \(\alpha\) make a constant angle with a fixed direction.
  Assume that \(\tau(s) \neq 0\), \( s \in I \), and prove that
  \begin{enumerate}[label=(\alph*)]
    \item \(\alpha\) is a helix if and only if \( \frac{k}{\tau} \) is constant.
    \item \(\alpha\) is a helix if and only if the lines containing \( n(s) \) and passing through \( \alpha(s) \) are parallel to a fixed plane.
    \item \(\alpha\) is a helix if and only if the lines containing \( b(s) \) and passing through \(\alpha(s)\) make a constant angle with a fixed direction.
    \item The curve
          \begin{equation}
          \label{eq:2}
          \alpha(s) = \left( \frac{a}{c} \int\sin\theta(s) ds, \frac{a}{c} \int\cos\theta(s) ds, \frac{b}{c} s \right),
          \end{equation}
          where \( c^2 = a^2 + b^2 \), is a helix, and that \( \frac{k}{\tau} = \frac{a}{b} \).
  \end{enumerate}
\end{problem}
\end{document}
